\mxunit{mada-t11-s-01}{Généralités sur les fonctions numériques}
\mxsubtitle{Analyse · première partie du programme de T11 S · fonctions numériques d'une variable réelle}
\begin{objectives}
\begin{itemize}
\item Lire une image et un antécédent sur une expression, sur un tableau et sur une courbe.
\item Déterminer l'ensemble de définition d'une fonction polynôme, rationnelle, irrationnelle ou contenant une valeur absolue.
\item Lire et démontrer un sens de variation à partir de la définition.
\item Établir qu'une fonction est paire, impaire ou ni l'une ni l'autre, et reconnaître un axe ou un centre de symétrie.
\item Établir qu'une fonction est périodique, trouver sa plus petite période pour $A\cos(ax+b)$ et $A\sin(ax+b)$, et en déduire un intervalle d'étude.
\end{itemize}
\end{objectives}
\begin{prerequisites}
Les fonctions de référence de seconde (affine, carré, cube, inverse, racine carrée, valeur absolue) ; résoudre une équation ou une inéquation du second degré ; dresser un tableau de signes ; le cercle trigonométrique, le radian, $\cos^2x+\sin^2x=1$ ; lire les coordonnées d'un point sur un graphique.
\end{prerequisites}
\begin{fluency}[5 min]
(a) Pour quels réels $x$ a-t-on $x+5\ge0$ ? (b) Résoudre $x^2=49$. (c) Calculer $\abs{-9}$ et $\abs{4-11}$. (d) Que vaut $\sqrt{(-6)^2}$ ? (e) Développer $(-x)^3$ et $(-x)^4$. (f) Donner $\cos(0)$, $\sin\bigl(\frac{\pi}{2}\bigr)$ et $\cos(2\pi)$.
\mxanswer{(a) $x\ge-5$ ; (b) $-7$ et $7$ ; (c) $9$ et $7$ ; (d) $6$ ; (e) $-x^3$ et $x^4$ ; (f) $1$, $1$, $1$.}
\end{fluency}

\section{Activité : deux manières de se répéter}
\begin{activity}[Un miroir et une frise]
Trois fonctions sont données par leurs valeurs aux entiers de $-3$ à $3$.
\begin{center}
\begin{tabular}{c|c|c|c|c|c|c|c}
$x$ & $-3$ & $-2$ & $-1$ & $0$ & $1$ & $2$ & $3$\\\hline
$f(x)$ & $7$ & $2$ & $-1$ & $-2$ & $-1$ & $2$ & $7$\\
$g(x)$ & $-15$ & $-6$ & $-1$ & $0$ & $1$ & $6$ & $15$\\
$h(x)$ & $2$ & $-1$ & $3$ & $2$ & $-1$ & $3$ & $2$
\end{tabular}
\end{center}
\textbf{1.} Pour chaque fonction, comparer la valeur en $x$ et la valeur en $-x$. Que remarque-t-on pour $f$ ? pour $g$ ? pour $h$ ?\par
\textbf{2.} On prolonge $h$ à tous les entiers en gardant le motif $2,-1,3$ qui se répète. Calculer $h(30)$ et $h(-7)$.\par
\textbf{3.} Peut-on affirmer, avec ce tableau seul, que $f(x)=f(-x)$ pour tous les réels $x$ ? Que faudrait-il de plus ?\mxlines{3}
\mxanswer{1. $f(-x)=f(x)$ pour les sept valeurs : la courbe semble symétrique par rapport à l'axe des ordonnées. $g(-x)=-g(x)$ : symétrie par rapport à l'origine. $h$ n'a aucune de ces deux propriétés ($h(-1)=3\ne h(1)=-1$) mais ses valeurs reviennent toutes les trois étapes : $h(x+3)=h(x)$. 2. $30=3\times10$, donc $h(30)=h(0)=2$ ; $-7=-9+2$, donc $h(-7)=h(2)=3$. 3. Non : sept valeurs ne disent rien des autres réels. Il faut une égalité $f(-x)=f(x)$ démontrée pour un $x$ quelconque de l'ensemble de définition.}
\end{activity}

\section{Fonction, image, antécédent}
\begin{definition}[Fonction numérique]
Une \textbf{fonction numérique} $f$ d'une variable réelle associe à chaque réel $x$ d'une partie $D$ de $\R$ \emph{un seul} réel, noté $f(x)$. La partie $D$ est l'\textbf{ensemble de définition} de $f$. Le réel $f(x)$ est l'\textbf{image} de $x$ ; si $f(x)=y$, on dit que $x$ est un \textbf{antécédent} de $y$.
\end{definition}
Un réel de $D$ a une seule image, mais un réel $y$ peut avoir plusieurs antécédents ou aucun : pour la fonction carré, $9$ en a deux ($-3$ et $3$) et $-1$ n'en a pas.
\begin{figurebox}[Lecture d'une image et d'antécédents sur la courbe de $f(x)=\frac{x^2}{3}-2$.]
\begin{tikzpicture}\begin{axis}[mxaxis,xmin=-5,xmax=5,ymin=-3,ymax=5,width=9cm,height=6.2cm,xtick={-4,-2,2,4},ytick={-2,2,4}]
\addplot[mxcurve,domain=-4.6:4.6,samples=100]{x*x/3-2};
\draw[mxaux](axis cs:3,0)--(axis cs:3,1);\draw[mxaux](axis cs:3,1)--(axis cs:0,1);
\draw[mxaux](axis cs:-3,1)--(axis cs:3,1);
\fill[mxpoint](axis cs:3,1)circle(2pt)node[above right]{$f(3)=1$};\fill[mxpoint](axis cs:-3,1)circle(2pt)node[above left]{$-3$};
\end{axis}\end{tikzpicture}
\end{figurebox}
Pour l'image de $3$ : on part de $3$ sur l'axe des abscisses, on monte jusqu'à la courbe, on lit l'ordonnée. Pour les antécédents de $1$ : on part de $1$ sur l'axe des ordonnées, on se déplace horizontalement et on relève l'abscisse de \emph{chaque} point d'intersection.
\begin{example}[Image et antécédents par calcul]
Soit $f(x)=x^2-4x+1$ sur $\R$. Image de $5$ : $f(5)=25-20+1=6$. Antécédents de $-2$ : $x^2-4x+1=-2\iff x^2-4x+3=0$ ; $\Delta=16-12=4$, d'où $x=\frac{4-2}{2}=1$ ou $x=\frac{4+2}{2}=3$. Antécédents de $-4$ : $x^2-4x+5=0$ a pour discriminant $16-20<0$ : aucun.
\end{example}

\section{Ensemble de définition}
Quand une fonction est donnée par une formule, son ensemble de définition est l'ensemble des réels pour lesquels le calcul de $f(x)$ est possible dans $\R$.
\begin{keypoint}[Les obstacles]
(1) Un quotient $\frac{A}{B}$ exige $B\ne0$. (2) Une racine $\sqrt{A}$ exige $A\ge0$. (3) Un quotient dont le dénominateur est une racine, $\frac{1}{\sqrt{A}}$, exige $A>0$. (4) Une valeur absolue $\abs{A}$ n'ajoute aucune condition. Un polynôme est défini sur $\R$ tout entier. La tangente exige $\cos x\ne0$.
\end{keypoint}
\begin{method}[Déterminer un ensemble de définition]
On repère chaque dénominateur et chaque racine ; on écrit une condition par obstacle ; on résout chaque condition (un tableau de signes aide pour le second degré ou un quotient) ; on garde les réels qui vérifient \emph{toutes} les conditions, donc l'intersection ; on écrit le résultat en intervalles.
\end{method}
\begin{example}[Deux obstacles]
Pour $g(x)=\dfrac{\sqrt{x+3}}{x-5}$ : sous la racine, $x+3\ge0$, soit $x\ge-3$ ; au dénominateur, $x\ne5$. Donc $D_g=\intfo{-3}{5}\cup\intoo{5}{+\infty}$. Le réel $-3$ est conservé (la racine vaut $0$), le réel $5$ est exclu.
\end{example}
\begin{example}[Un second degré sous la racine]
Pour $h(x)=\sqrt{x^2-x-12}$, on résout $x^2-x-12\ge0$. Les racines sont $-3$ et $4$ ($\Delta=49$) ; le trinôme est positif à l'extérieur : $D_h=\intof{-\infty}{-3}\cup\intfo{4}{+\infty}$.
\end{example}
\begin{pitfall}
Simplifier avant de chercher le domaine fait perdre des points exclus : $\frac{x^2-9}{x-3}$ n'est pas défini en $3$, même si l'expression simplifiée $x+3$ l'est. On détermine l'ensemble de définition sur la formule d'origine.
\end{pitfall}

\section{Sens de variation}
\begin{definition}[Croissance, décroissance]
Soit $f$ définie sur un intervalle $I$. Elle est \textbf{croissante} sur $I$ si, pour tous $a,b$ de $I$, $a\le b\Rightarrow f(a)\le f(b)$ ; \textbf{décroissante} si $a\le b\Rightarrow f(a)\ge f(b)$. Avec $a<b$ et des inégalités strictes, on parle de croissance ou de décroissance \textbf{stricte}. Elle est \textbf{monotone} sur $I$ si elle est croissante ou décroissante sur $I$.
\end{definition}
\begin{pitfall}
Un sens de variation se dit toujours sur un intervalle. La fonction inverse est décroissante sur $\intoo{-\infty}{0}$ et sur $\intoo{0}{+\infty}$, mais pas sur leur réunion : pour $a=-2<b=3$, on a $f(a)=-\frac12<f(b)=\frac13$, l'ordre n'est pas renversé. Il faut écrire deux intervalles, pas un seul.
\end{pitfall}
\begin{theorem}[Somme et multiple]
Soit $f$ et $g$ croissantes sur $I$ : alors $f+g$ est croissante sur $I$. Soit $f$ monotone sur $I$ et $\lambda\ne0$ : $\lambda f$ varie comme $f$ si $\lambda>0$ et en sens contraire si $\lambda<0$.
\end{theorem}
\begin{proof}
Pour $a\le b$ dans $I$, $f(a)\le f(b)$ et $g(a)\le g(b)$ ; en additionnant, $f(a)+g(a)\le f(b)+g(b)$. Si $f$ est croissante et $\lambda>0$, multiplier $f(a)\le f(b)$ par $\lambda$ conserve l'ordre ; si $\lambda<0$, cela le renverse, donc $\lambda f$ est décroissante. Les autres cas sont identiques.
\end{proof}
\begin{example}[Lire un tableau de variation]
$f(x)=-2x^2$ : le carré est décroissant sur $\intof{-\infty}{0}$ et croissant sur $\intfo{0}{+\infty}$ ; en multipliant par $-2<0$, $f$ est croissante sur $\intof{-\infty}{0}$ et décroissante sur $\intfo{0}{+\infty}$, avec le maximum $f(0)=0$.
\end{example}

\section{Parité}
Pour comparer $f(x)$ et $f(-x)$, il faut que $-x$ soit lui aussi dans l'ensemble de définition.
\begin{definition}[Ensemble symétrique, fonction paire, fonction impaire]
Une partie $D$ de $\R$ est \textbf{symétrique par rapport à $0$} si $x\in D\Rightarrow-x\in D$. Soit $f$ de domaine symétrique. Elle est \textbf{paire} si $f(-x)=f(x)$ pour tout $x$ de $D$, \textbf{impaire} si $f(-x)=-f(x)$ pour tout $x$ de $D$.
\end{definition}
\begin{theorem}[Lecture sur la courbe]
Dans un repère orthogonal, $f$ est paire si et seulement si sa courbe est symétrique par rapport à l'axe des ordonnées ; $f$ est impaire si et seulement si sa courbe est symétrique par rapport à l'origine.
\end{theorem}
\begin{proof}
Le symétrique du point $M\pt{x;f(x)}$ par rapport à l'axe des ordonnées est $M'\pt{-x;f(x)}$ ; il est sur la courbe exactement quand $f(-x)=f(x)$. Par rapport à l'origine, le symétrique est $M''\pt{-x;-f(x)}$, qui est sur la courbe exactement quand $f(-x)=-f(x)$.
\end{proof}
\begin{figurebox}[Une fonction paire ($\frac{x^2}{3}-2$, en trait plein) et une fonction impaire ($\frac{x^3-6x}{10}$, en pointillés).]
\begin{tikzpicture}\begin{axis}[mxaxis,xmin=-5,xmax=5,ymin=-4,ymax=4,width=9cm,height=6cm,xtick={-4,-2,2,4},ytick={-4,-2,2,4}]
\addplot[mxcurve,domain=-4.5:4.5,samples=100]{x*x/3-2};
\addplot[mxcurve,dashed,mxobjtwo,domain=-4.5:4.5,samples=100]{(x^3-6*x)/10};
\end{axis}\end{tikzpicture}
\end{figurebox}
\begin{method}[Étudier la parité]
1. Déterminer l'ensemble de définition et vérifier qu'il est symétrique par rapport à $0$ ; sinon, la fonction n'est ni paire ni impaire. 2. Calculer $f(-x)$ pour un $x$ quelconque. 3. Comparer à $f(x)$ et à $-f(x)$. Pour montrer qu'une fonction n'est ni paire ni impaire, un seul contre-exemple numérique suffit.
\end{method}
\begin{example}
$f(x)=x^4-5x^2+1$ : $D=\R$ et $f(-x)=x^4-5x^2+1=f(x)$ : paire. $g(x)=\dfrac{x}{x^2-4}$ : $D=\R\setminus\ens{-2;2}$, symétrique, et $g(-x)=\dfrac{-x}{x^2-4}=-g(x)$ : impaire. $k(x)=\sqrt{x}$ : $D=\intfo{0}{+\infty}$ n'est pas symétrique, donc $k$ n'est ni paire ni impaire. $m(x)=(x+1)^2$ : $m(1)=4$ et $m(-1)=0$ ne sont ni égaux ni opposés.
\end{example}
\begin{pitfall}
Le calcul $f(-x)=f(x)$ n'a de sens que si $D$ est symétrique : $\frac{1}{x-1}$ n'est pas définie en $1$ mais l'est en $-1$, donc elle n'est ni paire ni impaire, sans aucun calcul.
\end{pitfall}
\begin{theorem}[Autres symétries]
Soit $f$ définie sur une partie $D$ avec $x\in D\Rightarrow 2a-x\in D$. La droite $x=a$ est un axe de symétrie de la courbe si et seulement si $f(2a-x)=f(x)$ pour tout $x$ de $D$ ; le point $\Omega\pt{a;b}$ en est un centre de symétrie si et seulement si $f(2a-x)+f(x)=2b$ pour tout $x$ de $D$.
\end{theorem}
\begin{proof}
La symétrie d'axe $x=a$ envoie $M\pt{x;y}$ sur $\pt{2a-x;y}$, et la symétrie de centre $\Omega$ envoie $M$ sur $\pt{2a-x;2b-y}$. La courbe est globalement invariante exactement quand ces images sont encore sur la courbe, ce qui s'écrit $f(2a-x)=f(x)$, respectivement $f(2a-x)=2b-f(x)$.
\end{proof}
\begin{example}
$f(x)=x^2-6x+5$ : $f(3+h)=(3+h)^2-6(3+h)+5=h^2-4$, valeur qui ne change pas quand $h$ devient $-h$ : l'axe $x=3$ est un axe de symétrie. $g(x)=2+\frac{3}{x-1}$ : $g(1+h)+g(1-h)=4+\frac3h-\frac3h=4$ : le point $\pt{1;2}$ est un centre de symétrie.
\end{example}

\section{Périodicité}
\begin{definition}[Fonction périodique]
Soit $T>0$. Une fonction $f$ est \textbf{périodique de période $T$} si, pour tout $x$ de son ensemble de définition, $x+T$ et $x-T$ appartiennent à l'ensemble de définition et $f(x+T)=f(x)$. La \textbf{plus petite période} strictement positive, quand elle existe, s'appelle \emph{la} période de $f$.
\end{definition}
Si $T$ est une période, tout multiple entier non nul $kT$ en est une : $f(x+2T)=f(x+T)=f(x)$, et de même pour les suivants et pour $-T$. Le motif de la courbe sur un intervalle de longueur $T$ se recopie par translations de vecteur $kT\vec\imath$.
\begin{figurebox}[Courbe de $f(x)=2\cos(2x)$ : le motif sur $\intff{0}{\pi}$ se recopie, la période est $\pi$.]
\begin{tikzpicture}\begin{axis}[mxaxis,xmin=-4,xmax=7,ymin=-3,ymax=3,width=11cm,height=5cm,xtick={0,1.5708,3.1416,4.7124,6.2832},xticklabels={$0$,$\frac\pi2$,$\pi$,$\frac{3\pi}2$,$2\pi$},ytick={-2,2}]
\addplot[mxcurve,domain=-3.8:6.8,samples=200]{2*cos(deg(2*x))};
\draw[mxaux](axis cs:0,2.6)--(axis cs:3.1416,2.6);\node[above]at(axis cs:1.57,2.6){$T=\pi$};
\end{axis}\end{tikzpicture}
\end{figurebox}
\begin{theorem}[Période de $A\cos(ax+b)$ et $A\sin(ax+b)$]
Soit $A\ne0$, $a\ne0$ et $b$ réels. Les fonctions $x\mapsto A\cos(ax+b)$ et $x\mapsto A\sin(ax+b)$ sont définies sur $\R$, périodiques, de plus petite période $\dfrac{2\pi}{\abs{a}}$.
\end{theorem}
\begin{proof}
Avec $T=\frac{2\pi}{\abs a}$ : $\cos(a(x+T)+b)=\cos(ax+b\pm2\pi)=\cos(ax+b)$, donc $T$ est une période. Soit $T'>0$ une période de $g(x)=A\cos(ax+b)$ et $x_0$ un réel où $\cos(ax_0+b)=1$, c'est-à-dire $g(x_0)=A$ est la valeur extrême. Alors $g(x_0+T')=A$, donc $\cos(ax_0+aT'+b)=1$, d'où $aT'\in2\pi\Z$ et $T'=\frac{2\pi k}{\abs a}$ avec $k\ge1$ entier : $T'\ge T$. Le sinus s'obtient de la même façon, ou par $\sin u=\cos(u-\frac\pi2)$.
\end{proof}
\begin{method}[Réduire l'intervalle d'étude]
Si $f$ est périodique de période $T$, on l'étudie sur un intervalle de longueur $T$, par exemple $\intff{-T/2}{T/2}$ ; si elle est en plus paire ou impaire, sur $\intff{0}{T/2}$ ; puis on complète la courbe par symétrie et par translations de vecteurs $kT\vec\imath$.
\end{method}
\begin{example}
$f(x)=3\sin(4x-1)$ a pour période $\frac{2\pi}{4}=\frac\pi2$. $g(x)=\cos x+\cos 2x$ a pour période $2\pi$ (les deux termes en ont une) et est paire : on l'étudie sur $\intff{0}{\pi}$. $\tan(x+\pi)=\tan x$ : la tangente est $\pi$-périodique.
\end{example}
\begin{pitfall}
La somme de deux fonctions périodiques n'est pas toujours périodique : $\sin x+\sin(\sqrt2\,x)$ ne l'est pas, car les périodes $2\pi$ et $\frac{2\pi}{\sqrt2}$ n'ont aucun multiple commun. Pour $\cos x+\cos 2x$, on a $2\pi=1\times2\pi=2\times\pi$ : un multiple commun existe.
\end{pitfall}
\begin{beyond}[Pour aller plus loin : un exemple sans période]
Les fonctions $x\mapsto x^2$ et $x\mapsto\frac1x$ ne sont pas périodiques : une période $T$ imposerait $(x+T)^2=x^2$ pour tout $x$, donc $2xT+T^2=0$ pour tout $x$, ce qui oblige $T=0$.
\end{beyond}

\section{Exercices}
\subsection*{Niveau 1}
\begin{exercise}[Images et antécédents]\mxlabel{img}
Soit $f(x)=x^2-4x+1$ sur $\R$. (1) Calculer $f(0)$, $f(-2)$, $f(\frac12)$. (2) Trouver les antécédents de $-2$. (3) Le réel $-4$ a-t-il un antécédent ?
\mxlines{4}
\end{exercise}
\begin{exercise}[Image d'un quotient]\mxlabel{quot}
Soit $g(x)=\dfrac{2x+1}{x-3}$. (1) Calculer $g(0)$, $g(4)$, $g(-1)$. (2) Résoudre $g(x)=2$, puis $g(x)=1$. Que peut-on dire de l'antécédent éventuel de $2$ ?
\mxlines{4}
\end{exercise}
\begin{exercise}[Ensembles de définition]\mxlabel{dom1}
Déterminer l'ensemble de définition de chaque fonction.
\begin{enumerate}
\item $f(x)=2x^3-5x+1$
\item $f(x)=\dfrac{1}{x-7}$
\item $f(x)=\dfrac{3x}{x^2-16}$
\item $f(x)=\sqrt{3x+6}$
\item $f(x)=\sqrt{8-2x}$
\item $f(x)=\dfrac{1}{\sqrt{x+3}}$
\end{enumerate}
\mxlines{5}
\end{exercise}
\begin{exercise}[Deux fonctions qui se ressemblent]\mxlabel{dom2}
Déterminer l'ensemble de définition de $f(x)=\dfrac{x+2}{x^2+4}$ et de $g(x)=\dfrac{x+2}{x^2-4}$. Écrire $g(x)$ plus simplement. L'expression simplifiée a-t-elle le même ensemble de définition ?
\mxlines{4}
\end{exercise}
\begin{exercise}[Sens de variation de fonctions simples]\mxlabel{var1}
Donner, sans calcul compliqué, le sens de variation de : $x\mapsto-2x+5$ sur $\R$ ; $x\mapsto3x^2$ sur $\intfo{0}{+\infty}$ puis sur $\intof{-\infty}{0}$ ; $x\mapsto-x^2$ sur $\intfo{0}{+\infty}$.
\mxlines{4}
\end{exercise}
\begin{exercise}[Pair, impair ou ni l'un ni l'autre]\mxlabel{par1}
Étudier la parité de chaque fonction.
\begin{enumerate}
\item $f(x)=x^4-3x^2+2$
\item $f(x)=x^3+5x$
\item $f(x)=(x+1)^2$
\item $f(x)=\dfrac{x}{x^2-1}$
\item $f(x)=\abs{x}+x^2$
\item $f(x)=\sqrt{x}$
\item $f(x)=\dfrac{1}{x-1}$
\end{enumerate}
\mxlines{6}
\end{exercise}
\begin{exercise}[Valeurs d'une fonction périodique]\mxlabel{per1}
Une fonction $u$ de période $3$ vérifie $u(0)=1$, $u(1)=3$, $u(2)=2$. Calculer $u(31)$, $u(-4)$ et $u(300)$.
\mxlines{3}
\end{exercise}
\begin{exercise}[Période d'une fonction trigonométrique]\mxlabel{per2}
Donner une période de chacune : $f(x)=\sin(3x)$ ; $g(x)=\cos\bigl(\frac x2\bigr)$ ; $h(x)=5\cos\bigl(2x+\frac\pi3\bigr)$.
\mxlines{3}
\end{exercise}
\subsection*{Niveau 2}
\begin{exercise}[Sous une racine, au dénominateur]\mxlabel{dom3}
Déterminer l'ensemble de définition de : (1) $f(x)=\sqrt{x^2-6x+8}$ ; (2) $g(x)=\dfrac{\sqrt{5-x}}{x^2-9}$ ; (3) $h(x)=\dfrac{1}{\abs x-3}$ ; (4) $k(x)=\sqrt{\dfrac{x-1}{x+4}}$ (tableau de signes).
\mxlines{6}
\end{exercise}
\begin{exercise}[Variation par la définition]\mxlabel{var2}
Montrer, à partir de la définition, que la fonction inverse est décroissante sur $\intoo{-\infty}{0}$.
\mxlines{4}
\end{exercise}
\begin{exercise}[Encadrer une image]\mxlabel{var3}
Une fonction $f$ est croissante sur $\intff{-1}{6}$ avec $f(-1)=3$ et $f(6)=9$. Que peut-on dire de $f(2)$ ? Que ne peut-on pas affirmer ?
\mxlines{3}
\end{exercise}
\begin{exercise}[Somme et produit de fonctions croissantes]\mxlabel{var4}
(1) Démontrer que la somme de deux fonctions décroissantes sur $I$ est décroissante sur $I$. (2) Les fonctions $f(x)=x$ et $g(x)=x$ sont croissantes sur $\intoo{-\infty}{0}$ ; leur produit $fg$ l'est-il ?
\mxlines{5}
\end{exercise}
\begin{exercise}[Un domaine pas symétrique]\mxlabel{par2}
Justifier sans calcul que $f(x)=\sqrt{x+2}$ n'est ni paire ni impaire. Même question pour $g(x)=\dfrac{1}{x-5}$.
\mxlines{3}
\end{exercise}
\begin{exercise}[Axe de symétrie]\mxlabel{axe}
Montrer que la droite $x=3$ est un axe de symétrie de la courbe de $f(x)=x^2-6x+5$ et que la droite $x=-1$ est un axe de la courbe de $g(x)=\abs{x+1}$.
\mxlines{4}
\end{exercise}
\begin{exercise}[Centre de symétrie]\mxlabel{centre}
Montrer que $\Omega\pt{1;2}$ est un centre de symétrie de la courbe de $f(x)=\dfrac{2x+1}{x-1}$ (on pourra écrire $f(x)=2+\frac{3}{x-1}$).
\mxlines{4}
\end{exercise}
\begin{exercise}[Algèbre de la parité]\mxlabel{par3}
Soit $f$ paire et $g$ impaire, toutes deux de domaine $\R$. Étudier la parité de $f\times g$, de $f^2$ et de $f+g$ (donner un exemple pour $f+g$). Montrer que la seule fonction à la fois paire et impaire sur $\R$ est la fonction nulle.
\mxlines{6}
\end{exercise}
\begin{exercise}[Valeurs d'une fonction périodique]\mxlabel{per3}
$f$ est de période $5$ sur $\R$ et $f(x)=x^2$ sur $\intfo{0}{5}$. Calculer $f(12)$, $f(-1)$ et $f(\num{17,5})$.
\mxlines{4}
\end{exercise}
\begin{exercise}[Période et parité d'une fonction trigonométrique]\mxlabel{per4}
(1) Montrer que $\cos^2x$ a pour période $\pi$ et qu'elle est paire. (2) Montrer que $\sin x+\cos x$ a pour période $2\pi$. (3) À quel intervalle réduire l'étude de $f(x)=\cos x+\cos2x$ ?
\mxlines{5}
\end{exercise}
\subsection*{Niveau 3}
\begin{exercise}[Les multiples d'une période]\mxlabel{per5}
Montrer que si $T$ est une période de $f$ sur $\R$, alors $-T$ et $3T$ en sont aussi. Si $S$ et $T$ sont deux périodes, montrer que $S+T$ en est une.
\mxlines{4}
\end{exercise}
\begin{exercise}[Impaire et croissante]\mxlabel{imp}
Soit $f$ impaire sur $\R$, croissante sur $\intfo{0}{+\infty}$. Montrer que $f(0)=0$ et que $f$ est croissante sur $\intof{-\infty}{0}$.
\mxlines{5}
\end{exercise}
\begin{exercise}[Axe et période]\mxlabel{axper}
Soit $f$ de période $T$, paire sur $\R$. Montrer que la droite $x=\frac T2$ est un axe de symétrie de la courbe.
\mxlines{4}
\end{exercise}
\begin{beyond}[Hors programme : minimalité d'une période]
Montrer que $\pi$ est la plus petite période de $\cos^2x$ : écrire $\cos^2x=\frac{1+\cos2x}{2}$ puis utiliser le théorème de la plus petite période.
\mxlines{3}
\end{beyond}

\section{Problèmes}
\begin{exercise}[Problème : une courbe reconstituée]\mxlabel{recon}
Une fonction $f$ est paire, de période $4$ sur $\R$ ; sur $\intff{0}{2}$ elle vaut $f(x)=2x$ pour $x\le1$ et $f(x)=2$ pour $x\ge1$.
\begin{enumerate}
\item Calculer $f(-\num{0,5})$, $f(3)$, $f(10)$ et $f(\num{7,25})$.
\item Tracer la courbe sur $\intff{-4}{4}$.
\item Quelle est la valeur maximale de $f$ ? Où est-elle atteinte ?
\end{enumerate}
\mxlines{6}
\end{exercise}
\begin{exercise}[Problème : la grande roue]\mxlabel{roue}
La hauteur, en mètres, d'une nacelle au-dessus du sol, $t$ minutes après le départ, est $h(t)=6-5\cos\bigl(\frac{\pi t}{10}\bigr)$.
\begin{enumerate}
\item Quelle est la période de $h$ ? Que représente-t-elle ?
\item Calculer $h(0)$, $h(10)$, $h(35)$ et $h(50)$.
\item Montrer que $h(10+u)=h(10-u)$ pour tout $u$ et interpréter.
\item À quels instants de $\intff{0}{20}$ la nacelle est-elle à $6$ m ?
\end{enumerate}
\mxlines{6}
\end{exercise}

\begin{summary}[À retenir]
\begin{itemize}
\item Une fonction a une seule image par réel de son domaine, mais un réel peut avoir plusieurs antécédents ou aucun.
\item Domaine : $B\ne0$ pour $\frac AB$, $A\ge0$ pour $\sqrt A$, $A>0$ pour $\frac{1}{\sqrt A}$ ; on croise toutes les conditions.
\item Parité : domaine symétrique, puis $f(-x)=f(x)$ (paire, axe des ordonnées) ou $f(-x)=-f(x)$ (impaire, origine).
\item Axe $x=a$ : $f(2a-x)=f(x)$ ; centre $\pt{a;b}$ : $f(2a-x)+f(x)=2b$.
\item Période : $f(x+T)=f(x)$ ; la période de $A\cos(ax+b)$ et de $A\sin(ax+b)$ est $\frac{2\pi}{\abs a}$.
\item Parité et période réduisent l'étude à une partie de l'ensemble de définition, le reste se déduit par symétrie et translation.
\end{itemize}
\end{summary}
\ifmxteacher\mxinclude{corriges/c01}\fi
